% !TeX root = ../../金融经济学_王江_讲义及习题.tex
% 来源：王江《金融经济学》，印刷第 33—47 页（PDF 第 48—62 页）。
% 补充：金融经济学_第3章 Arrow-Debreu经济.pptx，第 1—30 页。
% 习题参考：史上最全最完整的金融经济学答案(王江版).pdf，第 8—12 页。

\zhchapter{Arrow-Debreu 经济}\label{chap:wj-three}
\begingroup
\raggedbottom
\renewcommand{\thesection}{\thechapter.\arabic{section}}
\renewcommand{\thesubsection}{\thesection.\arabic{subsection}}
\renewcommand{\thesubsubsection}{\thesubsection.\arabic{subsubsection}}
\ctexset{
 section={number=\thesection,aftername=\quad},
 subsection={number=\thesubsection,aftername=\quad},
 subsubsection={number=\thesubsubsection,aftername=\quad}
}
% 独立计数器保留教材定义、定理和例题号；标签绑定真实计数器。
\newcounter{wjThreeDefinition}
\newcounter{wjThreeTheorem}
\newcounter{wjThreeExample}
\renewcommand{\thewjThreeDefinition}{\thechapter.\arabic{wjThreeDefinition}}
\renewcommand{\thewjThreeTheorem}{\thechapter.\arabic{wjThreeTheorem}}
\renewcommand{\thewjThreeExample}{\thechapter.\arabic{wjThreeExample}}
\newcommand{\WJThreeDefinition}[1]{\refstepcounter{wjThreeDefinition}\label{#1}\textbf{定义 \thewjThreeDefinition}\quad}
\newcommand{\WJThreeTheorem}[1]{\refstepcounter{wjThreeTheorem}\label{#1}\textbf{定理 \thewjThreeTheorem}\quad}
\newcommand{\WJThreeExample}[1]{\refstepcounter{wjThreeExample}\label{#1}\textbf{例 \thewjThreeExample}\quad}

\ifdefined\WJHandoutKnowledge
本章考察由全部状态或有证券构成的市场。先说明状态价格与市场完全性，再从个人最优消费推导市场均衡，最后分析均衡配置的效率。支付矩阵仍以状态为行、证券为列；$\theta_k$ 表示证券净交易量，$h_k$ 表示交易后总持仓。

\section{Arrow-Debreu 证券市场}
\label{sec:wj-three-securities}
\begin{tcolorbox}[knowledge]
\WJThreeDefinition{def:wj-three-claim}状态 $\omega$ \textbf{或有要求权}（state-$\omega$ contingent claim）是在 1 期状态为 $\omega$ 时支付 $1$ 单位消费品、其他状态下支付 $0$ 的证券，也称状态 $\omega$ 的 Arrow-Debreu 证券。
\end{tcolorbox}
若以 $\omega'$ 表示实际实现的状态，其支付为
\[
 x(\omega')=\mathbf{1}_{\{\omega\}}(\omega')=
 \begin{cases}1,&\omega'=\omega,\\0,&\omega'\ne\omega.\end{cases}
\]
\begin{tcolorbox}[example]
\WJThreeExample{ex:wj-three-1}Lucas 树经济在 1 期有两个可能状态 $a,b$，两只状态或有证券的支付向量分别为
\[
 x_a=\begin{bmatrix}1\\0\end{bmatrix},\qquad
 x_b=\begin{bmatrix}0\\1\end{bmatrix}.
\]
\end{tcolorbox}
\begin{tcolorbox}[knowledge,title=Arrow-Debreu 市场与单位支付矩阵,breakable=false]
每个可能状态对应一只状态证券，共有 $\Omega$ 只。由它们的完全集合构成的市场称为 \textbf{Arrow-Debreu 证券市场}。按状态顺序排列证券时，支付矩阵为
\[
 X^{A\text{-}D}=I_\Omega.
\]
\end{tcolorbox}

其中第 $\omega$ 列只在第 $\omega$ 行取 $1$，直接对应状态 $\omega$ 的单位支付。

\section{状态价格}
\label{sec:wj-three-prices}
\begin{tcolorbox}[knowledge,title=状态价格]
状态 $\omega$ 的阿罗证券在 0 期的价格记为 $\phi_\omega$，以 0 期消费品计量，称为\textbf{状态价格}（state price）。状态价格向量为 $\phi=(\phi_1,\ldots,\phi_\Omega)^{\mathsf T}$。
\end{tcolorbox}

% 来源：教师 PPT 第 3 页的台风状态证券情境。
\begin{tcolorbox}[example,title=课堂例题：只在台风状态支付的证券]
设明天有晴天、下雨、台风、地震四个互斥且穷尽的可能状态，按此顺序排列。某证券只在台风发生时支付 $1$ 单位消费品，其他状态支付为零，其支付向量为 $(0,0,1,0)^{\mathsf T}$。今天购买一份该证券所需支付的消费品数量，就是台风状态的状态价格。
\end{tcolorbox}
\begin{tcolorbox}[knowledge,title=状态价格的正性,breakable=false]
只考虑正概率状态。在无摩擦、允许自由买卖且无套利时，
\begin{equation}
 \phi\gg0.\label{eq:wj-three-positive-price}
\end{equation}
\end{tcolorbox}

若某个 $\phi_\omega\le0$，买入一份该证券的当期净支出非正，未来支付非负且以正概率为正，构成套利。状态价格同时反映时间价值与相应状态下资源的稀缺性，其总和一般不为 $1$。在一价定律下，支付向量 $x$ 的复制成本为 $\phi^{\mathsf T}x$；实际概率则描述各状态发生的可能性。

\section{市场的完全性}
\subsection{消费计划的复制与完全市场}
在 Arrow-Debreu 市场中，给定有限的未来消费向量 $c_1=(c_{11},\ldots,c_{1\Omega})^{\mathsf T}$，取总持仓 $h=c_1$，便有
\[
 \underbrace{X^{A\text{-}D}h}_{\text{组合支付}}=I_\Omega c_1
 =\underbrace{c_1}_{\text{目标消费}},\qquad
 \underbrace{\phi^{\mathsf T}h}_{\text{复制成本}}=\sum_{\omega=1}^{\Omega}\phi_\omega c_{1\omega}.
\]
\begin{tcolorbox}[knowledge]
\WJThreeDefinition{def:wj-three-complete}如果任一有限消费计划都可以通过有限成本的可交易证券组合来融资，则称证券市场为\textbf{完全市场}（complete market），也称完备市场。
\end{tcolorbox}
有限状态、有限价格下，这要求证券支付能够生成全部状态支付。Arrow-Debreu 市场因此是完全市场。完全性说明支付可以复制；个人能够负担的消费还须满足其财富约束。

\begin{tcolorbox}[example]
\WJThreeExample{ex:wj-three-2}在两状态 Lucas 树经济中，市场上只有支付为 $(1,1)^{\mathsf T}$ 的无风险债券。若要复制消费 $(2,1)^{\mathsf T}$，债券持仓 $h$ 必须同时满足 $h=2$ 与 $h=1$，无解。因此市场不完全。
\end{tcolorbox}

\subsection{完全性的判别}
\label{sec:wj-three-rank}
\begin{tcolorbox}[knowledge,title=完全市场的行满秩判据,breakable=false]
设一般支付矩阵 $X\in\mathrm R^{\Omega\times N}$，允许任意实数持仓。市场完全当且仅当
\[
 \operatorname{Im}(X)=\mathrm R^\Omega
 \quad\Longleftrightarrow\quad \operatorname{rank}(X)=\Omega.
\]
\end{tcolorbox}

\textbf{（完全市场 $\Rightarrow$ 行满秩）}每个单位状态支付 $\mathbf e_\omega$ 都可复制，所以这些构成 $\mathrm R^\Omega$ 基底的向量均在 $X$ 的列空间中，列空间就是 $\mathrm R^\Omega$，其维数为 $\Omega$。

\textbf{（行满秩 $\Rightarrow$ 完全市场）}若列空间维数为 $\Omega$，则它等于 $\mathrm R^\Omega$，每个有限消费向量都能写成 $Xh$。证券价格与组合分量均有限时，其复制成本 $S^{\mathsf T}h$ 也有限，故市场完全。

证券数至少等于状态数是必要条件；证券之间的线性相关性决定这些证券提供了多少独立支付方向。

% 来源：教师 PPT 第 9 页的三组支付矩阵。
\begin{tcolorbox}[example,title=课堂例题：三组支付矩阵的完全性]
判断下列三个市场是否完全：
\[
 X_1=\begin{bmatrix}1&3\\2&4\end{bmatrix},\quad
 X_2=\begin{bmatrix}1&2\\2&4\end{bmatrix},\quad
 X_3=\begin{bmatrix}1&0&1\\0&1&1\\1&1&1\end{bmatrix}.
\]
\end{tcolorbox}
$\det X_1=-2\ne0$，故第一个市场完全。$X_2$ 的第二列是第一列的两倍，秩为 $1<2$，故第二个市场不完全。$\det X_3=-1\ne0$，故第三个市场完全；其中虽未直接列出全部阿罗证券，也能通过组合复制它们。两只证券面对三个状态时，秩最多为 $2$，市场必不完全。

在相同的状态空间中，完全市场都能生成全部状态支付。若偏好、禀赋及对应的支付定价相同，这些市场提供相同的消费机会；它们可以使用不同的证券作为支付空间的基底。

\subsection{市场不完全与金融创新}
% 来源：教师 PPT 第 10、23 页；将金融创新归入完全性主题。
\begin{tcolorbox}[example,title=课堂例题：保险合约为何可能缺失]
失业保险可能降低投保人寻找工作的激励；全额手机屏幕保险可能降低使用保护措施的意愿。健康状况较差的人也可能更愿意购买健康保险，而保险人难以准确识别其风险类型。分析这些因素如何影响保险合约的供给。
\end{tcolorbox}
\begin{tcolorbox}[knowledge,title=道德风险与逆向选择,breakable=false]
投保后行为发生变化属于\textbf{道德风险}；投保前不同风险类型的选择差异属于\textbf{逆向选择}。
\end{tcolorbox}

加上损失可验证性不足和合约执行成本，一些状态支付便难以成为可交易合约。大量独立风险有利于保险人分散损失，系统性风险则需要额外的风险承担能力和融资安排。

健康保险中，若保险人难以区分风险类型，可能按投保群体的平均风险收取统一保费。低风险者认为保费过高而退出后，留下的投保群体平均风险上升，保险人便可能进一步提高保费，促使更多低风险者退出。这种“保费上升—低风险者退出—平均风险继续上升”的反馈，可能使保险覆盖范围收缩，甚至使某些合约退出市场。

金融创新通过增加原有列空间之外的新支付向量，扩大可交易支付空间。例如，巨灾债券提供特定灾害状态的支付，通胀挂钩债券提供与购买力变化相关的支付，信用衍生品提供与违约事件相关的支付。其对完全性的贡献取决于新支付是否增加矩阵的秩。

\section{参与者的优化}
\subsection{禀赋财富与预算约束}
\label{sec:wj-three-budget}
\begin{tcolorbox}[knowledge,title=初始持仓、总持仓与净交易,breakable=false]
给定状态价格 $\phi\gg0$，参与者的禀赋为 $e=(e_0,e_1^{\mathsf T})^{\mathsf T}$。在阿罗证券中，复制未来禀赋的初始持仓为 $\bar\theta=e_1$。记交易后总持仓为 $h$，净买入量为 $\theta=h-\bar\theta$，则
\[
 c_1=e_1+\theta=h.
\]
\end{tcolorbox}

\begin{tcolorbox}[knowledge,title=禀赋的市场财富,breakable=false]
禀赋在当期的市场价值为
\begin{equation}
 w=\underbrace{e_0}_{\text{当期禀赋}}
   +\underbrace{\phi^{\mathsf T}e_1}_{\text{未来禀赋市值}}.\label{eq:wj-three-wealth}
\end{equation}

这称为总财富或禀赋的市场财富。
\end{tcolorbox}

\begin{tcolorbox}[knowledge,title=价值预算约束,breakable=false]
证券净交易支出与可负担消费的关系为
\begin{equation}
 \begin{aligned}
 c_0+\phi^{\mathsf T}\theta&\le e_0\\
 \Longleftrightarrow\quad c_0+\phi^{\mathsf T}c_1
 &\le e_0+\phi^{\mathsf T}e_1=w.
 \end{aligned}\label{eq:wj-three-budget}
\end{equation}
\end{tcolorbox}

预算不等式允许自由处置资源；在不满足性下，最优选择花光预算。将禀赋按市场价格计入财富后，个人可直接在消费空间中选择最优计划。

\begin{tcolorbox}[knowledge,title=扩展价格与消费优化,breakable=false]
记扩展价格向量 $\hat\phi=(1,\phi^{\mathsf T})^{\mathsf T}$，其中 $\phi_0=1$。参与者的问题为
\begin{equation}
 \begin{aligned}
 \max_{c\in\mathrm R_+^{1+\Omega}}\quad&U(c)\\
 \mathrm{s.t.}\quad&\hat\phi^{\mathsf T}c\le w=\hat\phi^{\mathsf T}e.
 \end{aligned}\label{eq:wj-three-optimization}
\end{equation}
\end{tcolorbox}

\subsection{最优消费的存在性}
\begin{tcolorbox}[knowledge]
\WJThreeTheorem{thm:wj-three-existence}设 $\hat\phi\gg0$、$0\le w<\infty$，且 $U$ 在 $\mathrm R_+^{1+\Omega}$ 上连续，则问题~\eqref{eq:wj-three-optimization}（预算约束下的效用最大化）有最优解。
\end{tcolorbox}
\textbf{证明。}预算集 $B=\{c\ge0:\hat\phi^{\mathsf T}c\le w\}$ 包含零向量，因而非空。每个分量满足
\[
 0\le c_i\le\frac{w}{\phi_i}\le\frac{w}{\min_{0\le j\le\Omega}\phi_j},
 \qquad i=0,\ldots,\Omega,
\]

所以预算集有界。若 $c^{(n)}\in B$ 且 $c^{(n)}\to c$，取极限仍有 $c\ge0$、$\hat\phi^{\mathsf T}c\le w$，故 $B$ 闭。在有限维空间中，非空有界闭集为紧集，连续效用函数在其上取得最大值。

对数效用的定义域为严格正消费集；例~\ref{ex:wj-three-3} 将结合其边界行为构造唯一最优解。

\subsection{最优消费的一阶条件}
\begin{tcolorbox}[knowledge,title=边际效用向量,breakable=false]
设 $U$ 在消费集的邻域上连续可微，且各项边际效用严格为正：
\begin{equation}
 DU(c)\gg0,\qquad
 DU(c)=(\partial_0U(c),\ldots,\partial_\Omega U(c))^{\mathsf T}.\label{eq:wj-three-marginal-utility}
\end{equation}

其中 $\partial_0U$ 表示当期消费的边际效用，$\partial_\omega U$ 表示未来状态 $\omega$ 消费的边际效用。
\end{tcolorbox}

% 教材印刷第 38 页及 PPT 第 14、17 页使用 +mu；按非负乘子约定改为 -mu。
% 凹性是本处充分性证明的附加条件，区别于前章偏好凸性（效用拟凹）。
\begin{tcolorbox}[knowledge]
\WJThreeTheorem{thm:wj-three-kkt}设 $w>0$、$\hat\phi\gg0$，效用函数 $U$ 在消费集的邻域上连续可微，并满足~\eqref{eq:wj-three-marginal-utility}（各项边际效用严格为正）。若 $c^*$ 是问题~\eqref{eq:wj-three-optimization}（预算约束下的效用最大化）的最优解，则存在 $\lambda>0$、$\mu_i\ge0$，使
\begin{subequations}
\begin{align}
 \partial_iU(c^*)&=\lambda\phi_i-\mu_i,\qquad i=0,\ldots,\Omega,\label{eq:wj-three-kkt-stationarity}\\
 \hat\phi^{\mathsf T}c^*&=w,\label{eq:wj-three-kkt-budget}\\
 \mu_i c_i^*&=0,\qquad i=0,\ldots,\Omega.\label{eq:wj-three-kkt-complementarity}
\end{align}
\end{subequations}
若进一步假设 $U$ 为凹函数，则可行消费 $c^*$ 为最优解，当且仅当它满足上述条件（一阶条件、预算用尽条件与互补松弛条件）。
\end{tcolorbox}

\textbf{（最优解 $\Rightarrow$ 一阶条件）}边际效用严格为正，故预算若有剩余便可增加消费，最优时必有~\eqref{eq:wj-three-kkt-budget}。又因 $w>0$，至少一个分量 $c_j^*>0$。记 $\mathbf e_i$ 为第 $i$ 个坐标方向的单位向量。对两个正消费分量 $i,j$，沿
\[
 d=\frac{\mathbf e_i}{\phi_i}-\frac{\mathbf e_j}{\phi_j}
\]

作充分小的正、反向调整均可行，且 $\hat\phi^{\mathsf T}d=0$。最优性要求 $DU(c^*)^{\mathsf T}d=0$，所以所有正消费分量的单位支出边际效用相同。记其共同值为 $\lambda=\partial_jU(c^*)/\phi_j>0$。

若 $c_i^*=0$，仅沿增加第 $i$ 项、减少正消费第 $j$ 项的方向作小调整可行。最优性要求该方向导数非正，即 $\partial_iU(c^*)/\phi_i\le\lambda$。于是令 $\mu_i=\lambda\phi_i-\partial_iU(c^*)\ge0$；正消费分量取 $\mu_i=0$，便同时得到~\eqref{eq:wj-three-kkt-stationarity} 与~\eqref{eq:wj-three-kkt-complementarity}。

\textbf{（一阶条件 $+$ 效用凹性 $\Rightarrow$ 最优解）}对任一可行消费 $a$，由凹函数的切平面不等式，
\[
 \begin{aligned}
 U(a)-U(c^*)&\le DU(c^*)^{\mathsf T}(a-c^*)\\
 &=\lambda\underbrace{\hat\phi^{\mathsf T}(a-c^*)}_{\le0}
 -\underbrace{\sum_{i=0}^{\Omega}\mu_i a_i}_{\ge0}
 +\underbrace{\sum_{i=0}^{\Omega}\mu_i c_i^*}_{=0}\le0.
 \end{aligned}
\]

因此 $c^*$ 至少与每个可行消费一样好，确为最优解。若 $U$ 严格凹，最优消费唯一。

这些条件也可由拉格朗日函数
\[
 \mathcal L=U(c)+\lambda(w-\hat\phi^{\mathsf T}c)+\sum_{i=0}^{\Omega}\mu_i c_i
\]

读出。$\lambda$ 是增加一单位财富的边际价值；$\mu_i$ 则对应消费非负约束。正消费分量的单位支出边际效用为 $\lambda$；零消费分量至多为 $\lambda$，此时继续减少该消费受到非负约束限制。

% 来源：教材印刷第 40 页；相对于非负消费集区分内点与边界。
\begin{tcolorbox}[knowledge,title=内点解与角点解,breakable=false]
若最优消费的所有分量都严格为正，称为\textbf{内点解}（interior solution）；若至少一个分量为零，则位于非负消费集的边界，通常称为\textbf{角点解}（corner solution）。
\end{tcolorbox}

\begin{tcolorbox}[example,title=课堂例题：角点消费与互补松弛]
设 $U(c_0,c_1)=2c_0+c_1$，两项价格均为 $1$，财富为 $1$。求最优消费及相应乘子。
\end{tcolorbox}
当期消费的单位支出边际效用为 $2$，高于未来消费的 $1$，故 $c^*=(1,0)$。取 $\lambda=2$、$\mu_0=0$、$\mu_1=1$，有 $\partial_1U=1=2-1$ 与 $\mu_1c_1^*=0$，满足全部条件。

\subsection{内点解与边际替代率}
\begin{tcolorbox}[knowledge,title=内点最优与边际替代率,breakable=false]
若最优消费 $c_k\gg0$，非负约束乘子均为零，由~\eqref{eq:wj-three-kkt-stationarity}（边际效用、价格与约束乘子的关系）得
\begin{equation}
 DU_k(c_k)=\lambda_k\hat\phi,\qquad
 \partial_0U_k=\lambda_k,\quad\partial_\omega U_k=\lambda_k\phi_\omega.\label{eq:wj-three-interior}
\end{equation}

因此跨期、跨状态消费的边际替代率满足
\begin{equation}
 \frac{\partial_\omega U_k(c_k)}{\partial_0U_k(c_k)}=\phi_\omega
 \quad\Longleftrightarrow\quad
 \partial_0U_k(c_k)=\frac1{\phi_\omega}\partial_\omega U_k(c_k).\label{eq:wj-three-mrs}
\end{equation}
\end{tcolorbox}

在预算内增加 $\mathrm dc_{1\omega}$ 单位状态 $\omega$ 消费，需要减少当期消费 $\mathrm dc_0=-\phi_\omega\mathrm dc_{1\omega}$，对应
\[
 \mathrm dU_k=
 \underbrace{(\partial_\omega U_k-\phi_\omega\partial_0U_k)}_{\text{预算内转移的边际净效用}}
 \mathrm dc_{1\omega}.
\]

内点最优时该系数为零，即市场交换比例与个人边际替代率一致。$\phi_\omega$ 表示为未来该状态增加一单位消费所放弃的当期消费量。

图~\ref{fig:wj-budget-tangency} 以两商品为例，给出预算线与无差异曲线相切的内点最优解。

% !TeX root = ../../../金融经济学_王江_讲义及习题.tex
% 教学补图：第3章内点条件。参考431计算讲义的预算线与无差异曲线风格。
\begin{center}
\begin{minipage}{\linewidth}
\centering
\begin{tikzpicture}[>=Stealth,font=\small,x=1.35cm,y=1.05cm]
 \fill[blue!5] (0,0)--(4,0)--(0,4)--cycle;
 \draw[->] (0,0)--(4.55,0) node[right] {$c_0$};
 \draw[->] (0,0)--(0,4.55) node[above] {$c_1$};
 \draw[red!70!black,thick] (0,4)--(4,0);
 \draw[blue!70!black,thick,domain=1:4,samples=70] plot (\x,{4/\x});
 \draw[blue!40!black,dashed,domain=.65:4,samples=60] plot (\x,{2/\x});
 \draw[dotted,gray] (2,0)--(2,2)--(0,2);
 \fill (2,2) circle(2pt) node[above right] {$c^*$};
 \node[below] at (4,0) {$w$}; \node[left] at (0,4) {$w/\phi$};
 \node[blue!70!black,right] at (4,1) {$U^*$};
 \node[red!70!black,align=left] at (3.5,3.25) {预算线斜率 $-1/\phi$};
\end{tikzpicture}
\par\smallskip\small
两商品示例的预算为 $c_0+\phi c_1=w$。
图取 $\phi=1,w=4,U=\log c_0+\log c_1$；切点满足 $\partial_1U/\partial_0U=\phi$。
\WJFigure{fig:wj-budget-tangency}{预算线与最优消费切点}
\end{minipage}
\end{center}


% 来源：教材印刷第 39—40 页，PPT 第 17—18 页；末项分母统一为 phi_b。
\begin{tcolorbox}[example]
\WJThreeExample{ex:wj-three-3} 考虑“Lucas 树经济”，1 期有两个概率相等的状态 $a$ 和 $b$。假设 1 期的两个可能状态的状态价格分别为 $\phi_{a}$ 和 $\phi_{b}$。考虑一个参与者，他的禀赋为 $(e_{0};e_{1a};e_{1b})$。他具有如下对数形式的效用函数：
\[
    U(c_{0},c_{1a},c_{1b})=\log c_{0}+\tfrac12(\log c_{1a}+\log c_{1b})
\]

他的最优消费/组合选择是什么？
\end{tcolorbox}
\begin{solution}
给定状态价格和他的禀赋，他的总财富是 $w=e_{0}+\phi_{a}e_{1a}+\phi_{b}e_{1b}$。他的最优化问题是
\[
    \begin{aligned}
        \max_{\{c_{0},c_{1a},c_{1b}\}}\ &\log c_{0}+\tfrac12(\log c_{1a}+\log c_{1b})\\
        \mathrm{s.t.}\quad &w-(c_{0}+\phi_{a}c_{1a}+\phi_{b}c_{1b})=0\\
        &c_{0},c_{1a},c_{1b}>0
    \end{aligned}
\]

假设 $w>0$、$\phi_a,\phi_b>0$。对数效用的消费域为 $c_0,c_{1a},c_{1b}>0$；最优解在内部，可直接使用等式预算的拉格朗日函数
\[
    \mathcal{L}=\log c_{0}+\tfrac12(\log c_{1a}+\log c_{1b})+\lambda\bigl(w-c_0-\phi_a c_{1a}-\phi_b c_{1b}\bigr).
\]

一阶条件为
\[
  \frac1{c_0}=\lambda,\qquad
  \frac1{2c_{1a}}=\lambda\phi_a,\qquad
  \frac1{2c_{1b}}=\lambda\phi_b,\qquad
  c_0+\phi_a c_{1a}+\phi_b c_{1b}=w.
\]

当消费趋近 $0$ 时，对数效用趋于负无穷、边际效用趋于正无穷，因此最优消费在每一时期、每一状态都严格为正。由上述一阶条件得到最优消费：
\[
    c_{0}^{*}=\frac{1}{\lambda},\qquad
    c_{1a}^{*}=\frac{1}{\lambda}\frac{1}{2\phi_{a}},\qquad
    c_{1b}^{*}=\frac{1}{\lambda}\frac{1}{2\phi_{b}}
\]

把 $c^{*}$ 代入预算约束，我们可以得到 $\lambda$ 的解：
\[
    \frac{1}{\lambda}+\frac{1}{2\lambda}+\frac{1}{2\lambda}=w\qquad\Rightarrow\qquad\lambda=\frac{2}{w}
\]

最后，我们有
\[
    c_{0}^{*}=\frac12 w,\qquad
    c_{1a}^{*}=\frac14\frac{w}{\phi_{a}},\qquad
    c_{1b}^{*}=\frac14\frac{w}{\phi_{b}}
\]

参与者将一半财富用于当期消费，另一半用于未来消费。给定财富 $w$，各未来状态的最优消费与状态价格成反比。价格变化还会通过 $w=e_0+\phi^{\mathsf T}e_1$ 改变禀赋财富，因此价格变动的总效应同时包含消费成本与财富的变化。

由于效用在正消费域上严格凹，上述可行的一阶条件解是唯一的最优解。各项最优支出为
\[
 \underbrace{c_0^*}_{w/2},\qquad
 \underbrace{\phi_a c_{1a}^*}_{w/4},\qquad
 \underbrace{\phi_b c_{1b}^*}_{w/4}.
\]

状态证券的初始持仓为 $\bar\theta=e_1$，交易后总持仓为 $h^*=c_1^*$，故净买入量为
\[
 \theta^*=h^*-\bar\theta
 =\begin{bmatrix}w/(4\phi_a)-e_{1a}\\w/(4\phi_b)-e_{1b}\end{bmatrix}.
\]

\end{solution}


\section{市场均衡}
\subsection{均衡的定义与市场出清}
\label{sec:wj-three-equilibrium}
\begin{tcolorbox}[knowledge,title=Arrow-Debreu 均衡,breakable=false]
均衡由一组状态价格 $\phi$ 和消费配置 $\{c_k\}_{k=1}^K$ 构成，满足：
\begin{enumerate}
 \item 每个参与者为价格接受者，在自己的预算集中选择最优消费：
 \[
 c_k\in\operatorname*{arg\,max}_{c\ge0,\ \hat\phi^{\mathsf T}c\le\hat\phi^{\mathsf T}e_k}U_k(c).
 \]
 \item 当期及未来每个状态的商品市场分别出清：
 \[
 \underbrace{\sum_k c_{k,0}=\sum_k e_{k,0}}_{\text{当期商品出清}},\qquad
 \underbrace{\sum_k c_{k,1\omega}=\sum_k e_{k,1\omega}}_{\text{状态 }\omega\text{ 商品出清}},
 \quad\omega=1,\ldots,\Omega.
 \]
\end{enumerate}
\end{tcolorbox}
这里的 $1+\Omega$ 种商品分别是当期消费品与未来各状态的或有消费品。上述定义正是\textbf{瓦尔拉斯均衡}（竞争均衡）在本模型中的形式。

\begin{tcolorbox}[knowledge,title=Walras 法则,breakable=false]
在不满足性下，各参与者花光预算。令 $z=\sum_k(c_k-e_k)$，加总预算等式即得 \textbf{Walras 法则}：
\[
 \hat\phi^{\mathsf T}z=\sum_k\hat\phi^{\mathsf T}(c_k-e_k)=0.
\]
\end{tcolorbox}

若价格严格为正，$1+\Omega$ 个市场中任意 $\Omega$ 个已出清，剩下一个也出清。因此求解价格可使用其中 $\Omega$ 个条件，最后一式用于核验。未来出清条件逐状态成立。

\subsection{均衡的存在性}
% 教材定理 3.3（第 42 页）及 PPT 第 22 页未单列生存性条件。
% 保留原命题归属，同时给出足以使用有限维交换经济存在定理的一组条件。
\begin{tcolorbox}[knowledge]
\WJThreeTheorem{thm:wj-three-equilibrium-existence}在两期、有限状态的 Arrow-Debreu 交换经济中，若每人的偏好连续、凸且严格单调（即公理~\ref{axiom:wj-two-1} 所述的不满足性），并且每人的禀赋向量均严格为正，则存在竞争均衡；均衡价格各分量严格为正，可将当期价格归一为 $1$。
\end{tcolorbox}
这是有限商品交换经济的均衡存在结论：把每个日期、状态下的消费看作一种商品，完全的状态证券将个人预算化为该交换经济的价值预算。其一般证明需要不动点方法，本章使用该存在结论。

对于含零禀赋分量的具体经济，可直接检验最优消费与市场出清，例~\ref{ex:wj-three-4} 即采用这种方法。
% 原书定理 3.3 未逐项给出生存性条件。此处保留严格正个人禀赋作为明确的充分条件。
% 仅有非负禀赋不足以保证严格正价格均衡：禀赋 (1,0)、效用 sqrt(c0)+sqrt(c1) 时未来总资源为零，而任意严格正价格下的最优未来消费均为正。


\subsection{均衡中的边际替代率}
\subsubsection{共同边际估值与互利交易}
\begin{tcolorbox}[knowledge,title=均衡中的共同边际替代率,breakable=false]
在效用可微、边际效用严格为正且均衡消费均为内点时，每个人都满足~\eqref{eq:wj-three-mrs}（内点最优时，边际替代率等于状态价格），因而
\begin{equation}
 \frac{\partial_\omega U_k(c_k)}{\partial_0U_k(c_k)}
 =\phi_\omega
 =\frac{\partial_\omega U_{k'}(c_{k'})}{\partial_0U_{k'}(c_{k'})},
 \qquad\forall k,k',\omega.\label{eq:wj-three-common-mrs}
\end{equation}
\end{tcolorbox}

共同的状态价格使不同参与者关于同一状态消费的边际估值一致。

% 来源：教材印刷第 42 页、PPT 第 22 页；展开双方交易的可行性与效用变化。
记参与者 $k$ 对状态 $\omega$ 消费的边际估值为
\[
 m_{k,\omega}=\frac{\partial_\omega U_k(c_k)}{\partial_0U_k(c_k)}.
\]

考虑两个消费均为内点的参与者，假设 $m_{k,\omega}>\phi_\omega>m_{k',\omega}$。让 $k$ 从 $k'$ 买入 $t>0$ 份状态 $\omega$ 证券，支付 $\phi_\omega t$ 单位当期消费品。双方消费变化分别为
\[
 \begin{aligned}
 (\Delta c_{k,0},\Delta c_{k,1\omega})&=(-\phi_\omega t,t),\\
 (\Delta c_{k',0},\Delta c_{k',1\omega})&=(\phi_\omega t,-t).
 \end{aligned}
\]

其余状态消费保持不变。双方各自的消费变化按市场价格计值为零，并且逐日期、逐状态相互抵消，所以交易既满足个人预算，也保持总资源可行。内点条件保证 $t$ 足够小时，双方消费仍非负。沿此交易方向，
\[
 \begin{aligned}
 \left.\frac{\mathrm dU_k}{\mathrm dt}\right|_{t=0}
 &=\partial_0U_k\underbrace{(m_{k,\omega}-\phi_\omega)}_{>0}>0,\\
 \left.\frac{\mathrm dU_{k'}}{\mathrm dt}\right|_{t=0}
 &=\partial_0U_{k'}\underbrace{(\phi_\omega-m_{k',\omega})}_{>0}>0.
 \end{aligned}
\]

因此存在足够小的正交易量，使双方效用都严格提高。这样的配置尚未达到个人最优与市场均衡。若两人的边际估值不同，任何介于二者之间的交易价格都提供相应的互利交换机会；均衡时共同状态价格与各人的边际估值一致，这类进一步交易动机随之消失。

\subsubsection{两状态经济的均衡}
% 来源：教材印刷第 42—43 页，PPT 第 24 页。
\begin{tcolorbox}[example]
\WJThreeExample{ex:wj-three-4} 考虑一个经济，在 1 期有两个概率相等的状态 $a$ 和 $b$。经济中有参与者 1 和 2，他们具有的禀赋分别为：
\begin{center}
\WJEventTree[left={$e_{1}{:}100$}]{0,0}
\qquad
\WJEventTree[left={$e_{2}{:}0$}]{200,50}
\end{center}
两个参与者都具有如下形式的对数效用函数：
\[
    U(c)=\log c_{0}+\tfrac12(\log c_{a}+\log c_{b})
\]

求均衡状态价格、消费配置和证券净交易量。
\end{tcolorbox}
\begin{solution}
在市场上存在一组完全的状态或有证券可以交易。因为有两个状态，因而只有两个状态或有证券。

现在我们开始分析这个经济的均衡。从给定交易证券价格下参与者的最优化问题开始。记 $\phi=\begin{bmatrix}\phi_{a} \\ \phi_{b}\end{bmatrix}$ 为状态价格（向量），即两个状态或有证券的价格。我们可以定义每个参与者的财富为 $w=\hat{\phi}^{\top}e$，这里 $\hat{\phi}=\begin{bmatrix}1 \\ \phi\end{bmatrix}$，而 $e$ 是他的禀赋。这时，最优化问题变成了：
\[
    \begin{aligned}
        \max_{c}\ &\log c_{0}+\tfrac12(\log c_{a}+\log c_{b})\\
        \mathrm{s.t.}\quad &c_{0}+\phi_{a}c_{a}+\phi_{b}c_{b}=w
    \end{aligned}
\]

由例~\ref{ex:wj-three-3}，最优消费为
\[
    c_{k,0}^{*}=\tfrac12 w_{k},\qquad
    c_{k,a}^{*}=\tfrac14\frac{w_{k}}{\phi_{a}},\qquad
    c_{k,b}^{*}=\tfrac14\frac{w_{k}}{\phi_{b}},\qquad k=1,2
\]

在状态价格 $\phi_{a},\phi_{b}$ 下，参与者按禀赋市值计的财富为 $w_{k}=e_{k,0}+\phi_{a}e_{k,a}+\phi_{b}e_{k,b}$（0 期消费品价格为 $1$）。代入禀赋 $(100;0,0)$ 与 $(0;200,50)$：
\[
    w_{1}=100+0\cdot\phi_{a}+0\cdot\phi_{b}=100,\qquad
    w_{2}=0+200\phi_{a}+50\phi_{b}=200\phi_{a}+50\phi_{b}.
\]

\textbf{0 期商品市场出清}（两人 0 期消费之和等于经济 0 期总禀赋 $100$）。由例~\ref{ex:wj-three-3}，$c_{k,0}^{*}=\tfrac12 w_{k}$，故
\[
    c_{1,0}^{*}+c_{2,0}^{*}=\tfrac12 w_{1}+\tfrac12 w_{2}=\tfrac12(w_{1}+w_{2})=100
    \quad\Longrightarrow\quad w_{1}+w_{2}=200.
\]

把 $w_{1},w_{2}$ 代入 $w_{1}+w_{2}=200$：
\[
    100+(200\phi_{a}+50\phi_{b})=200.
\]

移项得 $200\phi_{a}+50\phi_{b}=100$；两边同除以 $50$，得到与状态价格有关的第一个方程
\[
    4\phi_{a}+\phi_{b}=2.
\]

\textbf{1 期·状态 $a$ 商品市场出清}（状态 $a$ 下总消费 $200$）。由 $c_{k,a}^{*}=\tfrac14 w_{k}/\phi_{a}$，
\[
    c_{1,a}^{*}+c_{2,a}^{*}=\tfrac{w_{1}}{4\phi_{a}}+\tfrac{w_{2}}{4\phi_{a}}
    =\frac{w_{1}+w_{2}}{4\phi_{a}}=200
    \quad\Longrightarrow\quad w_{1}+w_{2}=800\phi_{a}.
\]

与 $w_{1}+w_{2}=200$ 联立得 $\phi_{a}=\dfrac14$，从而 $\phi_{b}=1$。状态 $b$ 出清 $\dfrac{w_{1}+w_{2}}{4\phi_{b}}=50$ 自动成立。此时 $w_{1}=w_{2}=100$，最优消费为
\[
    \begin{aligned}
        c_{1,0}^{*}=c_{2,0}^{*}&=50\\
        c_{1,a}^{*}=c_{2,a}^{*}&=100\\
        c_{1,b}^{*}=c_{2,b}^{*}&=25
    \end{aligned}
\]

将消费代回边际替代率，可逐状态检验均衡的一阶条件。对 $k=1,2$，都有
\[
 \begin{aligned}
 m_{k,a}&=\frac{1/(2c_{k,a}^*)}{1/c_{k,0}^*}
 =\frac{50}{2\times100}=\frac14=\phi_a,\\
 m_{k,b}&=\frac{1/(2c_{k,b}^*)}{1/c_{k,0}^*}
 =\frac{50}{2\times25}=1=\phi_b.
 \end{aligned}
\]

两人在每个状态的边际估值相同，并分别等于该状态的市场价格。

初始状态证券持仓分别为 $(0,0)^{\mathsf T}$ 与 $(200,50)^{\mathsf T}$，因此
\[
 \theta_1^*=\begin{bmatrix}100\\25\end{bmatrix},\qquad
 \theta_2^*=\begin{bmatrix}-100\\-25\end{bmatrix},\qquad
 \underbrace{\theta_1^*+\theta_2^*=0}_{\text{证券净交易出清}}.
\]

参与者 1 支出 $100/4+25=50$，参与者 2 获得同额收入，二人的当期消费均为 $50$。状态 $b$ 的总禀赋较少，每单位消费更稀缺，因此即使两状态概率相同，仍有 $\phi_b>\phi_a$。

\end{solution}


\subsubsection{证券缺失与交易成本}
% 来源：教师 PPT 第 22 页关于不完全市场的讨论。
若一般证券市场的支付矩阵为 $X$、证券价格为 $S$，在内点、无交易限制下，对净交易量求一阶条件只得到
\[
 X^{\mathsf T}m_k=S,\qquad
 (m_k)_\omega=\frac{\partial_\omega U_k}{\partial_0U_k}.
\]

两人之差满足 $X^{\mathsf T}(m_k-m_{k'})=0$。当 $X$ 行满秩时，其转置的零空间只有零向量，故边际替代率向量相同；市场不完全时，未被证券支付覆盖的方向仍可能存在估值差异。这说明完全性在消除进一步交易动机中的作用。

% 来源：教材印刷第 44 页、PPT 第 22 页；用买卖价差说明交易成本的影响。
\textbf{交易成本与无交易区间。}即使某状态证券可以交易，买卖价差也会改变上述条件。设买入一份须付 $\phi_\omega^{\mathrm{ask}}$，卖出一份只能得到 $\phi_\omega^{\mathrm{bid}}$，且
\[
 0<\phi_\omega^{\mathrm{bid}}\le\phi_\omega^{\mathrm{ask}}.
\]

仅考虑这种单位交易成本，且没有持仓限制。对消费为内点的参与者，增加少量买入带来的边际效用为 $\partial_0U_k(m_{k,\omega}-\phi_\omega^{\mathrm{ask}})$，增加少量卖出带来的边际效用为 $\partial_0U_k(\phi_\omega^{\mathrm{bid}}-m_{k,\omega})$。因此两种方向均无严格的一阶改善，需要
\[
 \underbrace{\phi_\omega^{\mathrm{bid}}\le m_{k,\omega}\le\phi_\omega^{\mathrm{ask}}}_{\text{边际上的无交易区间}}.
\]

不同参与者的边际估值可以位于这个区间的不同位置，而都没有增加交易的边际动机。若效用为凹函数，这些方向条件也足以保证单独交易该证券无法提高效用。无价差时，区间收缩为一个价格，便恢复~\eqref{eq:wj-three-common-mrs} 的等式条件。

% 来源：教师 PPT 第 28 页；将共同边际估值与可复制证券定价相连。
\subsubsection{状态价格与证券定价}
在完全的状态证券市场中，证券 $n$ 在未来各状态的支付为 $x_n=(x_{1n},\ldots,x_{\Omega n})^{\mathsf T}$。持有 $x_{\omega n}$ 份状态 $\omega$ 的阿罗证券，便能复制其支付，因此在无摩擦、一价定律下，
\[
 S_n=\underbrace{\phi^{\mathsf T}x_n}_{\text{复制组合的成本}}
 =\sum_{\omega=1}^{\Omega}\phi_\omega x_{\omega n}.
\]

均衡条件~\eqref{eq:wj-three-common-mrs}（各参与者的边际替代率共同等于状态价格）进一步把这里的状态价格与各参与者的边际替代率联系起来。资产价格由各状态支付按状态价格加权求和得到，而状态价格综合反映偏好、状态概率与资源稀缺性。

\subsection{均衡配置的效率}
\label{sec:wj-three-welfare}
可行配置逐日期、逐状态满足 $\sum_k c_k=\sum_k e_k$。在同一价格下计价，有
\begin{equation}
 \hat\phi^{\mathsf T}\sum_k c_k=\hat\phi^{\mathsf T}\sum_k e_k.\label{eq:wj-three-feasibility-value}
\end{equation}
\begin{tcolorbox}[knowledge]
\WJThreeTheorem{thm:wj-three-welfare}在本章无摩擦、完全的状态证券市场中，若偏好满足不满足性（公理~\ref{axiom:wj-two-1}，喜多厌少），则竞争均衡配置是 \textbf{\hyperref[def:wj-two-9]{Pareto 最优的}}：不存在另一可行配置，使所有人至少一样好、至少一人严格更好。
\end{tcolorbox}
\textbf{证明。}记均衡价格为 $\hat\phi\gg0$，均衡配置为 $\{c_k\}$，$w_k=\hat\phi^{\mathsf T}e_k=\hat\phi^{\mathsf T}c_k$。反设存在 Pareto 改善的可行配置 $\{a_k\}$。

先证每个弱改善者均有 $\hat\phi^{\mathsf T}a_k\ge w_k$。若其成本小于 $w_k$，令剩余预算为 $b=w_k-\hat\phi^{\mathsf T}a_k>0$，取
\[
 d=\frac{b}{2\sum_{i=0}^{\Omega}\phi_i}(1,\ldots,1)^{\mathsf T}\gg0.
\]
则 $a_k+d$ 的成本为 $w_k-b/2<w_k$，仍在预算内。由不满足性，$a_k+d\succ_k a_k\succeq_k c_k$，与 $c_k$ 最优矛盾。

对严格改善者 $j$，必有 $\hat\phi^{\mathsf T}a_j>w_j$；否则 $a_j$ 本来就在预算内，却严格优于 $c_j$，同样矛盾。因此
\[
 \underbrace{\hat\phi^{\mathsf T}\sum_k a_k}_{\text{改善配置的总成本}}
 >\sum_k w_k
 =\underbrace{\hat\phi^{\mathsf T}\sum_k e_k}_{\text{总禀赋价值}},
\]
与可行性 $\sum_k a_k=\sum_k e_k$ 矛盾。故均衡配置 Pareto 最优。

定理~\ref{thm:wj-three-welfare} 即\textbf{福利经济学第一定理}。它说明分散的个人选择可以在完全、无摩擦的竞争市场中达到跨期、跨状态配置效率。这一证明直接比较可行配置的预算成本，适用于包括角点消费在内的均衡配置。

\begin{tcolorbox}[knowledge,title=效率与分配]
Pareto 效率排除了使所有人不受损、至少一人受益的进一步调整，但有效配置仍可对应不同的福利分配。第二福利定理研究在附加条件下，如何通过一次性财富再分配支持给定的有效配置；这一问题将在完全市场的资源配置章节进一步讨论。
\end{tcolorbox}

\section*{本章小结}
\addcontentsline{toc}{section}{本章小结}
\begin{tcolorbox}[summary]
\begin{enumerate}
 \item 每个状态的阿罗证券提供该状态下的一单位支付，全部阿罗证券构成的支付矩阵为单位阵。
 \item 无套利时状态价格严格为正；它是或有消费的价格，实际概率则衡量状态发生的可能性。
 \item 完全市场能够复制全部状态支付。有限维、可自由组合时，完全性等价于支付矩阵行满秩。
 \item 给定价格与财富，个人先在预算集内优化；内点最优时，边际替代率等于状态价格。
 \item 均衡同时满足个人最优与逐状态市场出清。在相应存在条件下，完全市场具有竞争均衡；在可微、内点条件下，各人的边际替代率一致。
 \item 第一福利定理说明均衡配置 Pareto 最优，将状态证券的交易与资源配置效率联系起来。
\end{enumerate}
\end{tcolorbox}
\clearpage
\fi

\section*{课后习题}
\addcontentsline{toc}{section}{课后习题}
% !TeX root = ../../金融经济学_王江_讲义及习题.tex
% 教材第 3 章练习 1—4，印刷第 45—47 页（PDF 第 60—62 页）。教材例 3.3、3.4 位于本章正文。

\newcounter{wjThreeExercise}
\renewcommand{\thewjThreeExercise}{\thechapter.\arabic{wjThreeExercise}}
\newcommand{\WJThreeProblem}[2]{\refstepcounter{wjThreeExercise}\label{#1}\problemno{\thewjThreeExercise#2}}

\par\noindent
\begin{minipage}{\linewidth}

\begin{problem}
\WJThreeProblem{prob:wj-three-1}{} 考虑如下经济，在 1 期有两个可能状态 $a$ 和 $b$：
\begin{center}
\WJEventTree{a,b}
\end{center}

\begin{enumerate}[label=(\alph*)]
    \item 描述所有 Arrow-Debreu 证券的支付向量。记这些证券的价格向量为 $\phi$。
    \item 考虑一个拥有如下禀赋的参与者：
    \begin{center}
    \WJEventTree[left={$0$}]{2,1}
    \end{center}
    把他的禀赋表示成 Arrow-Debreu 证券的组合。
    \item 计算他的金融财富。写出他的预算集。
    \item 假设参与者的效用函数如下：
    \[
        U(c_{0},c_{1a},c_{1b})=-e^{-c_{0}}-\tfrac12\bigl(e^{-c_{1a}}+e^{-c_{1b}}\bigr).
    \]

    不考虑消费的非负约束，写出他的优化问题。求解他的最优消费选择。
    \item 讨论他的消费如何依赖于 Arrow-Debreu 证券的价格向量 $\phi$。
    \item 证明在某些价格下，他（在某些时期/状态下）的消费可能是负的。
\end{enumerate}
\end{problem}
\end{minipage}\par
\begin{solution}
\begin{enumerate}[label=(\alph*)]
    \item Arrow-Debreu 证券的支付向量是 $X_{a}=\begin{bmatrix}1 \\ 0\end{bmatrix}$、$X_{b}=\begin{bmatrix}0 \\ 1\end{bmatrix}$，对应的价格向量为 $\phi=(\phi_a,\phi_b)^{\mathsf T}$。
    \item $2X_{a}+X_{b}$；
    \item 0 期禀赋为 $0$，1 期禀赋是 $2$ 单位状态 $a$ 或有要求与 $1$ 单位状态 $b$ 或有要求，金融财富为 $w=2\phi_{a}+\phi_{b}$。在消费集 $\mathrm{R}_{+}^{3}$ 上，预算集是
    \[
        \bigl\{c\in\mathrm{R}_{+}^{3}:c_{0}+\phi_{a}c_{1a}+\phi_{b}c_{1b}\le w\bigr\}.
    \]
    \item 本问允许 $c\in\mathrm{R}^3$，并设 $\phi_a>0$、$\phi_b>0$。若预算有剩余，增加 $c_0$ 即可提高效用，故最优时预算取等号。优化问题为
    \[
        \begin{aligned}
            \max_{\{c_{0},c_{1a},c_{1b}\}}\ &\bigl\{-e^{-c_{0}}-\tfrac12\bigl(e^{-c_{1a}}+e^{-c_{1b}}\bigr)\bigr\}\\
            \mathrm{s.t.}\quad &c_{0}+\phi_{a}c_{1a}+\phi_{b}c_{1b}=w=2\phi_{a}+\phi_{b}
        \end{aligned}
    \]

    拉格朗日函数为
    \[
        \mathcal{L}=-e^{-c_{0}}-\tfrac12\bigl(e^{-c_{1a}}+e^{-c_{1b}}\bigr)+\lambda\bigl(w-c_{0}-\phi_{a}c_{1a}-\phi_{b}c_{1b}\bigr)
    \]

    对消费求导，一阶条件为
    \[
        e^{-c_{0}}=\lambda,\qquad
        \tfrac12 e^{-c_{1a}}=\lambda\phi_{a},\qquad
        \tfrac12 e^{-c_{1b}}=\lambda\phi_{b}
    \]

    由第一式得 $\lambda>0$。对三式分别取对数，得到
    \[
        c_{0}^{*}=-\log\lambda,\qquad
        c_{1a}^{*}=-\log(2\lambda\phi_{a}),\qquad
        c_{1b}^{*}=-\log(2\lambda\phi_{b})
    \]

    代入预算约束：
    \[
        -\log\lambda-\phi_{a}\log(2\lambda\phi_{a})-\phi_{b}\log(2\lambda\phi_{b})=w
    \]
    即 $-(1+\phi_{a}+\phi_{b})\log\lambda-\phi_{a}\log(2\phi_{a})-\phi_{b}\log(2\phi_{b})=w$。记 $S=1+\phi_{a}+\phi_{b}$，则
    \[
        \log\lambda=-\frac{w+\phi_{a}\log(2\phi_{a})+\phi_{b}\log(2\phi_{b})}{S}
    \]

    从而 $c_{0}^{*}=-\log\lambda$，且 $c_{1a}^{*}=-\log\lambda-\log(2\phi_{a})$、$c_{1b}^{*}=-\log\lambda-\log(2\phi_{b})$。代入 $w=2\phi_{a}+\phi_{b}$ 得到
    \[
        \begin{aligned}
            c_{0}^{*}&=\frac{2\phi_{a}+\phi_{b}+\phi_{a}\log(2\phi_{a})+\phi_{b}\log(2\phi_{b})}{1+\phi_{a}+\phi_{b}}\\
            c_{1a}^{*}&=\frac{2\phi_{a}+\phi_{b}+\phi_{b}\log(2\phi_{b})-\log(2\phi_{a})-\phi_{b}\log(2\phi_{a})}{1+\phi_{a}+\phi_{b}}\\
            c_{1b}^{*}&=\frac{2\phi_{a}+\phi_{b}+\phi_{a}\log(2\phi_{a})-\log(2\phi_{b})-\phi_{a}\log(2\phi_{b})}{1+\phi_{a}+\phi_{b}}
        \end{aligned}
    \]

    该计划满足预算等式与一阶条件。效用函数严格凹，预算等式给出的可行集是凸集，故这是唯一的全局最优解。
    \item 设状态 $j$ 的价格 $\phi_j$ 变动，另一状态的价格 $\phi_k$ 固定。由 (d) 的一阶条件 $c_{1j}^{*}=c_0^{*}-\log(2\phi_j)$ 和 $c_{1k}^{*}=c_0^{*}-\log(2\phi_k)$ 得
    \[
        \frac{\partial c_{1k}^{*}}{\partial\phi_j}
        =\frac{\partial c_0^{*}}{\partial\phi_j},
        \qquad
        \frac{\partial c_{1j}^{*}}{\partial\phi_j}
        =\frac{\partial c_0^{*}}{\partial\phi_j}-\frac{1}{\phi_j}.
    \]

    对预算等式 $c_0^*+\phi_a c_{1a}^*+\phi_b c_{1b}^*=2\phi_a+\phi_b$ 关于 $\phi_j$ 求导，再用上面两式将未来消费的导数写成当期消费的导数。沿用 (d) 的 $S=1+\phi_a+\phi_b$，得到
    \[
        S\frac{\partial c_0^*}{\partial\phi_j}
        +\underbrace{c_{1j}^*}_{\text{原消费的支出增量}}
        -\underbrace{1}_{\text{相对消费减少所节省的支出}}
        =\underbrace{e_{1j}}_{\text{禀赋增值}},
        \qquad e_{1a}=2,\quad e_{1b}=1.
    \]

    当期消费与价格未变的状态消费同向变化。$\phi_a$ 上升时，$c_0^*$ 与 $c_{1b}^*$ 的变化方向由 $(3-c_{1a}^*)/S$ 的符号决定。$\phi_b$ 上升时，$c_0^*$ 与 $c_{1a}^*$ 的变化方向由 $(2-c_{1b}^*)/S$ 的符号决定。价格上升的状态还要从当期消费的导数中扣去 $1/\phi_j$，差值为正时，该状态的消费才会上升。

    例如，固定 $\phi_b$，小幅提高 $\phi_a$。在下表的两组价格处，偏导数的正负分别表示对应的消费量上升或下降。
    \begin{center}
        \begin{tabular}{@{}crrr@{}}
            $(\phi_a,\phi_b)$ &
            $\partial c_0^*/\partial\phi_a$ &
            $\partial c_{1a}^*/\partial\phi_a$ &
            $\partial c_{1b}^*/\partial\phi_a$ \\
            \hline
            $(0.1,0.1)$ & $1.17$ & $-8.83$ & $1.17$ \\
            $(0.01,1)$  & $-0.87$ & $-100.87$ & $-0.87$
        \end{tabular}
    \end{center}
    \item 当 $\phi_{a}=\phi_{b}=0.1$ 时，$w=0.3$，$S=1.2$，
    \[
        c_{0}^{*}=\frac{0.3+0.2\log(0.2)}{1.2}\approx -0.0182<0.
    \]

    该计划满足预算等式，但 $c_{0}^{*}<0$，因而不属于 (c) 的预算集。
\end{enumerate}
\end{solution}

\par\noindent
\begin{minipage}{\linewidth}
\begin{problem}
\WJThreeProblem{prob:wj-three-2}{} 考虑一个在 1 期只有一个可能状态的经济。（在这种情况下不存在不确定性。）参与者 1 的 0 期禀赋为 100 而 1 期禀赋为 1，即他的禀赋向量为 $\begin{bmatrix}100 \\ 1\end{bmatrix}$。他的偏好可以表示成如下形式：
\[
    U(c_{0},c_{1})=\log c_{0}+\rho\log c_{1}.
\]

有一只可交易证券，它的 0 期价格为 1、1 期支付为 $1+r_{F}$。这里，$r_{F}$ 是利率。
\begin{enumerate}[label=(\alph*)]
    \item 假设利率 $r_{F}$ 是给定的，导出参与者对证券的需求。
    \item 假设参与者 1 是经济中的唯一参与者。描述市场出清条件。
    \item 求解均衡利率。
    \item 均衡利率如何依赖于偏好参数 $\rho$？解释所得到的结论。
\end{enumerate}
\end{problem}
\end{minipage}\par
\begin{solution}
\begin{enumerate}[label=(\alph*)]
    \item 设 $\rho>0$。每份债券当期花费 $1$，$r_F$ 是扣除本金后的净利率，$R=1+r_F>0$ 是含本金的到期支付。参与者 1 净买入 $\theta_1$ 份债券，当期支出 $\theta_1$，到期收到 $R\theta_1$，因此
    \[
        c_{1,0}=100-\theta_1,\qquad c_{1,1}=1+R\theta_1.
    \]

    将两期消费代入效用函数后，只需选择 $\theta_1$。对数效用要求 $c_{1,0}>0$、$c_{1,1}>0$，分别给出储蓄和借款的上限。于是有
    \[
        \begin{aligned}
            \max_{\theta_1}\quad &\log(100-\theta_1)+\rho\log(1+R\theta_1)\\
            \text{s.t.}\quad &-\frac{1}{R}<\theta_1<100.
        \end{aligned}
    \]

    对任意可行的 $\theta_1$，由前面两式确定的消费都满足预算等式，因此可以直接对 $\theta_1$ 求导，无需另写拉格朗日函数。用链式法则求导，内点最优时有
    \[
        \frac{\mathrm d}{\mathrm d\theta_1}
        \bigl[\log(100-\theta_1)+\rho\log(1+R\theta_1)\bigr]
        =-\frac{1}{100-\theta_1}+\frac{\rho R}{1+R\theta_1}
        =0\qquad\text{（一阶条件）}
    \]

    移项并交叉相乘得 $1+R\theta_1=\rho R(100-\theta_1)$。整理得 $(1+\rho)R\theta_1=100\rho R-1$，所以证券需求为
    \[
        \theta_1^{*}(R)=\frac{100\rho R-1}{(1+\rho)R}.
    \]

    代回两期预算式，得到
    \[
        c_{1,0}^{*}=\frac{100R+1}{(1+\rho)R},\qquad
        c_{1,1}^{*}=\frac{\rho(100R+1)}{1+\rho}.
    \]

    两期消费均为正，所以解在可行区间内。目标函数严格凹，因此这是一阶条件给出的唯一最优解。
    \item 出清要求所有参与者的最优证券净买入量之和为零。经济中只有参与者 1，因此出清条件为 $\theta_1^{*}(R)=0$。
    \item 由 (b)，均衡时 $\theta_1^*(R^*)=0$，其中 $R^*=1+r_F^*$。代入 (a) 求得的证券需求式，得到
    \[
        \theta_1^{*}(R^*)=\frac{100\rho R^*-1}{(1+\rho)R^*}=0.
    \]

    在 $\rho>0$、$R^*>0$ 的范围内，分母严格为正，故只需令分子为零。
    \[
        100\rho R^*-1=0
        \quad\Longrightarrow\quad R^*=\frac{1}{100\rho}
        \quad\Longrightarrow\quad r_F^{*}=\frac{1}{100\rho}-1.
    \]

    这个解满足 $R^*>0$。由于均衡时净买入量为零，两期消费就是初始禀赋 $(100,1)$。
    \item 均衡利率 $r_F^{*}=1/(100\rho)-1$ 随 $\rho$ 增大而下降。一份债券当期花费 $1$、未来支付 $R$，所以一单位未来消费的当期价格是 $q=1/R$。均衡时
    \[
      q^{*}=100\rho,\qquad R^{*}=\frac1{100\rho}.
    \]
    未来消费的相对权重提高，使参与者在给定利率下更愿意储蓄；均衡时单位未来支付的价格上升，相应的利率下降。
\end{enumerate}
\end{solution}

\par\noindent
\begin{minipage}{\linewidth}
\begin{problem}
\WJThreeProblem{prob:wj-three-3}{} 继续考虑练习~\ref{prob:wj-three-2} 中定义的经济。现在假设还有另外一个参与者 2，他的效用函数和参与者 1 的一样但他的禀赋为 $\begin{bmatrix}1 \\ 100\end{bmatrix}$（参与者 1 是$\begin{bmatrix}100 \\ 1\end{bmatrix}$）。
\begin{enumerate}[label=(\alph*)]
    \item 求解参与者 2 的最优消费/储蓄选择。
    \item 比较参与者 1 和 2 的最优消费/储蓄选择。
    \item 求解有两个参与者时的市场均衡。
    \item 计算均衡利率。它与前一情形（经济中只有参与者 1）下的均衡利率有何不同？
    \item 证明均衡配置是 Pareto 最优的。
\end{enumerate}
\end{problem}
\end{minipage}\par
\begin{solution}
练习~\ref{prob:wj-three-2} 只有参与者 1，出清时他的证券净买入量为零。本题的第 1 期仍只有一个可能状态。加入参与者 2 后，先分别求两人的证券需求，再用两人的净买入量之和为零求均衡利率。

\begin{enumerate}[label=(\alph*)]
    \item 沿用练习~\ref{prob:wj-three-2}(a) 的记号。参与者 2 净买入 $\theta_2$ 单位债券后，$c_{2,0}=1-\theta_2$，$c_{2,1}=100+R\theta_2$。把两期消费代入效用函数，并要求两者都为正，得到
    \[
        \begin{aligned}
            \max_{\theta_2}\quad &\log(1-\theta_2)+\rho\log(100+R\theta_2)\\
            \text{s.t.}\quad &-\frac{100}{R}<\theta_2<1.
        \end{aligned}
    \]

    用同样的链式法则对 $\theta_2$ 求导，内点最优时有
    \[
        \frac{\mathrm d}{\mathrm d\theta_2}
        \bigl[\log(1-\theta_2)+\rho\log(100+R\theta_2)\bigr]
        =-\frac{1}{1-\theta_2}+\frac{\rho R}{100+R\theta_2}
        =0\qquad\text{（一阶条件）}
    \]

    移项并交叉相乘得 $100+R\theta_2=\rho R(1-\theta_2)$，所以
    \[
        \theta_2^*(R)=\frac{\rho R-100}{(1+\rho)R}.
    \]

    将最优证券净买入量 $\theta_2^*(R)$ 代入两期消费式，得到
    \[
        c_{2,0}^{*}=1-\theta_2^*(R)=\frac{R+100}{(1+\rho)R},\qquad
        c_{2,1}^{*}=100+R\theta_2^*(R)=\frac{\rho(R+100)}{1+\rho}.
    \]

    由此还可发现，$c_{2,0}^*=\frac{1}{\rho}\cdot\frac{c_{2,1}^*}{R}=\frac{w_2}{1+\rho}$，其中 $w_2=1+100/R$ 是禀赋现值。

    两期消费都为正，且目标函数严格凹，所以这是唯一最优解。
    \item 两人的禀赋现值分别为 $w_1=100+1/R$、$w_2=1+100/R$。由两人的最优消费式可见，他们都把财富的 $1/(1+\rho)$ 用于当期消费，剩余财富用于购买未来消费。按当期价格计，两人用于当期消费和未来消费的支出之比都是 $1:\rho$，两项支出分别为
    \[
      c_{k,0}^*=\frac{w_k}{1+\rho},\qquad
      \underbrace{\frac{c_{k,1}^*}{R}}_{\text{未来消费的现值}}
      =\frac{\rho w_k}{1+\rho},\qquad k=1,2.
    \]

    给定 $R$，两期最优消费分别是财富 $w_k$ 的 $1/(1+\rho)$ 倍和 $\rho R/(1+\rho)$ 倍。这两个系数都为正，且对两人相同，所以财富较高者两期消费都较高。两人的财富之差为
    \[
      w_1-w_2=99\left(1-\frac1R\right).
    \]

    当 $R>1$ 时，$w_1-w_2>0$，所以参与者 1 的两期消费均较高。当 $R<1$ 时，$w_1-w_2<0$，所以参与者 2 的两期消费均较高。$R=1$ 时两人财富相等，两期消费也相同。

    再看储蓄和借款。两人的当期消费式为 $c_{1,0}=100-\theta_1$ 和 $c_{2,0}=1-\theta_2$。把最优当期消费代入这两式，解得净买入量
    \[
      \theta_1^*=100-c_{1,0}^*=\frac{100\rho R-1}{(1+\rho)R},\qquad
      \theta_2^*=1-c_{2,0}^*=\frac{\rho R-100}{(1+\rho)R}.
    \]

    净买入量为正表示储蓄，为负表示借款。两式的分母都为正，故
    \[
      \theta_1^*>0\iff R>\frac1{100\rho},\qquad
      \theta_2^*>0\iff R>\frac{100}{\rho}.
    \]

    当 $R$ 低于两个临界值时，两人都借款。高于两个临界值时，两人都储蓄。当 $1/(100\rho)<R<100/\rho$ 时，参与者 1 储蓄，参与者 2 借款。$R$ 恰好等于某人的临界值时，此人不交易。

    \item 均衡时，两人的证券净买入量相加为零。将练习~\ref{prob:wj-three-2}(a) 和本题 (a) 的需求代入出清条件，得到
    \[
        \underbrace{\frac{100\rho R^*-1}{(1+\rho)R^*}}_{\text{参与者 1 的证券净买入量}}
        +\underbrace{\frac{\rho R^*-100}{(1+\rho)R^*}}_{\text{参与者 2 的证券净买入量}}
        =\frac{101\rho R^*-101}{(1+\rho)R^*}=0.
    \]

    分母为正，令分子为零得 $R^*=1/\rho>0$。将 $R^*$ 代回两人的需求，得到
    \[
        \theta_1^*(R^*)=\frac{99\rho}{1+\rho},\qquad
        \theta_2^*(R^*)=-\frac{99\rho}{1+\rho}.
    \]

    参与者 1 储蓄，参与者 2 借款。由 (b)，最优消费满足 $c_{k,1}^*/R=\rho c_{k,0}^*$。代入均衡条件 $\rho R^*=1$，得到
    \[
        c_{k,1}^*=\rho R^*c_{k,0}^*=c_{k,0}^*,\qquad k=1,2.
    \]

    从各自的当期禀赋中扣除净买入量，再用两期消费相等的关系，得到
    \[
        c_{1,0}^*=c_{1,1}^*=\frac{100+\rho}{1+\rho},
        \qquad
        c_{2,0}^*=c_{2,1}^*=\frac{1+100\rho}{1+\rho}.
    \]

    两人各自选择了最优消费，证券净交易之和为零，两期总消费也都等于总禀赋 $101$。
    \item 有两位参与者时，参与者 1 的净买入量正好抵消参与者 2 的净卖出量。(c) 用出清条件 $\theta_1^*(R^*)+\theta_2^*(R^*)=0$ 求得 $R^*=1/\rho$。扣除本金 $1$，得到两人经济的均衡利率
    \[
      r_F^*=R^*-1=\frac{1}{\rho}-1.
    \]

    只有参与者 1 时，出清要求他自己的净买入量为零。由练习~\ref{prob:wj-three-2}(c)，此时的均衡利率为 $1/(100\rho)-1$。同一 $\rho>0$ 下，两人经济的均衡利率高出
    \[
      \left(\frac1\rho-1\right)-\left(\frac1{100\rho}-1\right)
      =\frac{99}{100\rho}>0.
    \]
    加入参与者 2 后，总禀赋从 $(100,1)$ 变为 $(101,101)$，未来商品相对更充裕。一单位未来商品的当期价格 $1/R^*$ 从 $100\rho$ 降为 $\rho$，因此均衡利率更高。

    \Needspace{13\baselineskip}
    \item 记 $c_k^*$ 为参与者 $k$ 的均衡消费计划。由 (c)，两期消费均严格为正，且 $c_{k,1}^*=\rho R^*c_{k,0}^*$。因此，未来消费相对于当期消费的跨期边际替代率为
    \[
        \frac{\partial_1U_k(c_k^*)}{\partial_0U_k(c_k^*)}
        =\rho\frac{c_{k,0}^*}{c_{k,1}^*}
        =\frac{1}{R^*}=\rho.
    \]

    这个比率衡量在效用不变时，参与者愿为增加未来消费放弃多少当期消费。在均衡点，它等于一单位未来商品的当期价格 $1/R^*$\ifdefined\WJHandoutKnowledge（见正文式~\eqref{eq:wj-three-mrs} 的内点最优条件）\fi。

    假设存在另一组可行消费计划 $\{a_k\}$，使两人的效用都不下降，且至少一人严格提高。记 $\Delta c_{k,t}=a_{k,t}-c_{k,t}^*$，其中 $t=0,1$。沿从 $c_k^*$ 到 $a_k$ 的线段，效用是凹的，所以在 $c_k^*$ 处用边际效用算出的线性变化，是实际效用增量的上界。也就是
    \[
      0\le U_k(a_k)-U_k(c_k^*)
      \le \partial_0 U_k(c_k^*)\,\Delta c_{k,0}
          +\partial_1 U_k(c_k^*)\,\Delta c_{k,1}.
    \]

    两边除以正数 $\partial_0 U_k(c_k^*)$，就把效用增量的上界换算成当期商品。利用上面的边际替代率，得到
    \[
      0\le \frac{U_k(a_k)-U_k(c_k^*)}{\partial_0 U_k(c_k^*)}
      \le \Delta c_{k,0}+\frac{\Delta c_{k,1}}{R^*}.
    \]

    右边正是按均衡价格计算的消费支出变化。两人的支出都不会减少，效用严格提高的人还必须增加支出，因此总支出会增加。但两期总资源都固定，任何可行的重新分配都满足 $\sum_k\Delta c_{k,0}=\sum_k\Delta c_{k,1}=0$，按同一组价格计算的总支出变化只能为零。两者矛盾，故均衡配置是 Pareto 最优的\ifdefined\WJHandoutKnowledge。这与正文第一福利定理（定理~\ref{thm:wj-three-welfare}）的计价思路相同\fi。\qedhere
\end{enumerate}
\end{solution}

\par\noindent
\begin{minipage}{\linewidth}
\begin{problem}
\WJThreeProblem{prob:wj-three-4}{} 在 1 期，经济有三个等可能的状态 $a$、$b$ 和 $c$。在证券市场上交易的 Arrow-Debreu 证券集合是完全的。经济中有两个参与者，参与者 1 和 2。他们的禀赋如下：
\begin{center}
\WJEventTree[left={$e_{1}{:}100$}]{0,0,0}
\qquad
\WJEventTree[left={$e_{2}{:}0$}]{200,100,50}
\end{center}
他们的偏好相同，形式如下：
\[
    U(c)=-\frac{1}{c_{0}}-\frac13\Bigl(\frac{1}{c_{1a}}+\frac{1}{c_{1b}}+\frac{1}{c_{1c}}\Bigr).
\]
\begin{enumerate}[label=(\alph*)]
    \item 求解每一参与者的最优消费/投资选择。
    \item 求解均衡状态价格。
    \item 证明对两个参与者的消费都有禀赋越高、消费越高的性质。
    \item 证明每一参与者的消费在总消费中所占的份额 $c_{k,1\omega}/e_{1\omega}$ 在所有状态下都是一样的，其中 $k=1,2$，而 $e_{1\omega}$ 是在状态 $\omega=a,b,c$ 下的 1 期总消费/禀赋。
\end{enumerate}
\end{problem}
\end{minipage}\par
\begin{solution}
练习~\ref{prob:wj-three-2} 只有一位参与者，练习~\ref{prob:wj-three-3} 增加到两位，两题的第 1 期都只有一个可能状态。本题仍有两位参与者，第 1 期有三个等可能的状态 $a,b,c$。前两题用一个利率描述跨期交换，本题要分别求三个状态下未来消费的价格 $\phi_a,\phi_b,\phi_c$。

\begin{enumerate}[label=(\alph*)]
    \item 先在给定的正状态价格下求个人需求\ifdefined\WJHandoutKnowledge（状态价格的正性见正文式~\eqref{eq:wj-three-positive-price}）\fi。消费不能为负，效用式的分母也不能为零，所以各项消费都必须为正。在这一范围内，效用严格递增、严格凹。按当期商品计价，禀赋财富为 $w_k=e_{k,0}+\sum_{\omega=a,b,c}\phi_\omega e_{k,1\omega}$\ifdefined\WJHandoutKnowledge（见正文式~\eqref{eq:wj-three-wealth}）\fi。代入两人的禀赋，得
    \[
       w_1=100,\qquad w_2=200\phi_a+100\phi_b+50\phi_c.
    \]

    按价值预算约束\ifdefined\WJHandoutKnowledge（正文式~\eqref{eq:wj-three-budget}）\fi，当期消费加上未来各状态消费的现值不能超过 $w_k$。效用严格递增，最优时用尽预算，因此两人的问题可以写为
    \[
      \begin{aligned}
       \max_{c_{k,0},c_{k,1\omega}>0}\quad&
       -\frac1{c_{k,0}}-\frac13\sum_{\omega=a,b,c}\frac1{c_{k,1\omega}}\\
       \mathrm{s.t.}\quad&c_{k,0}+\sum_{\omega=a,b,c}\phi_\omega c_{k,1\omega}=w_k.
      \end{aligned}
    \]

    以 $\lambda_k$ 为预算等式的乘子，写出拉格朗日函数
    \[
       \mathcal L_k=
       -\frac1{c_{k,0}}-\frac13\sum_{\omega=a,b,c}\frac1{c_{k,1\omega}}
       +\lambda_k\left(w_k-c_{k,0}-\sum_{\omega=a,b,c}\phi_\omega c_{k,1\omega}\right).
    \]

    对 $\lambda_k$ 求导并令导数为零，得到原预算等式。分别对 $c_{k,0}$ 和 $c_{k,1\omega}$ 求导并令导数为零，得到
    \[
       \frac1{c_{k,0}^2}=\lambda_k,\qquad
       \frac1{3c_{k,1\omega}^2}=\lambda_k\phi_\omega
       \qquad\text{（一阶条件）}.
    \]

    用第一式消去 $\lambda_k$，由消费为正取正平方根，得到 $c_{k,1\omega}=c_{k,0}/\sqrt{3\phi_\omega}$。代回预算等式，得
    \[
       w_k=c_{k,0}+\sum_{\omega=a,b,c}\phi_\omega\frac{c_{k,0}}{\sqrt{3\phi_\omega}}
       =c_{k,0}\underbrace{\left(1+\sum_{\omega=a,b,c}\sqrt{\frac{\phi_\omega}{3}}\right)}_{A}.
    \]

    因此两人的最优消费为
    \[
       c_{k,0}^*=\frac{w_k}{A},\qquad
       c_{k,1\omega}^*=\frac{w_k}{A\sqrt{3\phi_\omega}},
       \quad k=1,2,\ \omega=a,b,c.
    \]

    这些消费均严格为正且满足预算，目标函数严格凹，所以这是一阶条件给出的唯一最优解。每份阿罗证券在对应状态支付一单位商品，因此参与者在状态 $\omega$ 的最优净买入量为 $\theta_{k,\omega}^*=c_{k,1\omega}^*-e_{k,1\omega}$\ifdefined\WJHandoutKnowledge（见正文第~\ref{sec:wj-three-budget}~节的净交易关系）\fi。参与者 1 在三个状态都没有未来禀赋，所需的未来消费都由证券买入取得。参与者 2 则分别从三个状态的最优消费中扣除禀赋 $200,100,50$，得到净买入量。

    \item 记题中的 1 期总禀赋为 $E_{1\omega}=\sum_k e_{k,1\omega}$。当期出清要求 $\sum_k c_{k,0}^*=100$。由 (a) 的最优消费关系，在均衡状态价格 $\phi_\omega^*$ 下，
    \[
        \sum_k c_{k,1\omega}^*
        =\frac{\sum_k c_{k,0}^*}{\sqrt{3\phi_\omega^*}}
        =\frac{100}{\sqrt{3\phi_\omega^*}}.
    \]

    令各状态的总消费等于该状态的总禀赋，得到
    \[
       \frac{100}{\sqrt{3\phi_a^*}}=200,\qquad
       \frac{100}{\sqrt{3\phi_b^*}}=100,\qquad
       \frac{100}{\sqrt{3\phi_c^*}}=50.
    \]

    由 $\sqrt{3\phi_\omega^*}=100/E_{1\omega}$，平方后除以 $3$，得 $\phi_\omega^*=\tfrac13(100/E_{1\omega})^2$。代入三个状态的总禀赋，得到
    \[
       \phi_a^*=\frac1{12},\qquad\phi_b^*=\frac13,\qquad\phi_c^*=\frac43.
    \]

    三个状态价格都为正。$\phi_c^*=4/3$ 表示在 0 期要用 $4/3$ 单位商品，才能买到状态 $c$ 下的一单位消费。状态 $c$ 发生的概率是题目给定的 $1/3$。

    将均衡价格代入 (a) 的 $A$ 和禀赋财富，得
    \[
       A^*=1+\frac16+\frac13+\frac23=\frac{13}{6},
       \qquad
       w_2^*=\frac{200}{12}+\frac{100}{3}+\frac43\cdot50=\frac{350}{3}.
    \]

    例如，参与者 1 的当期消费为 $100/(13/6)=600/13$。由于 $\sqrt{3\phi_a^*}=1/2$，他在状态 $a$ 的消费为 $2c_{1,0}^*=1200/13$。用 (a) 的公式算得各项消费如下。
    \begin{center}
    \begin{tabular}{lcccc}
      \toprule
      & 0 期 & 1 期状态 $a$ & 1 期状态 $b$ & 1 期状态 $c$ \\
      \midrule
      参与者 1 & $600/13$ & $1200/13$ & $600/13$ & $300/13$ \\
      参与者 2 & $700/13$ & $1400/13$ & $700/13$ & $350/13$ \\
      总消费 & $100$ & $200$ & $100$ & $50$ \\
      \bottomrule
    \end{tabular}
    \end{center}

    各列的总消费都等于对应的总禀赋。按 (a) 的净交易公式，参与者 1 在三个状态的净买入量就是表中的未来消费量。参与者 2 在每个状态的净卖出量恰好与之相等，因此三个状态的证券市场也都出清。

    \Needspace{8\baselineskip}
    \item 这里比较的是各状态的总禀赋。由 (a) 的个人最优消费式和 (b) 的状态出清式，对任一参与者 $k$ 都有
    \[
       c_{k,1\omega}^*
       =\frac{c_{k,0}^*}{\sqrt{3\phi_\omega^*}}
       =\frac{c_{k,0}^*}{100}\,E_{1\omega}.
    \]

    同一个人的系数 $c_{k,0}^*/100$ 为正，且在三个状态都相同。由于 $E_{1a}>E_{1b}>E_{1c}$，两人各自都有 $c_{k,1a}^*>c_{k,1b}^*>c_{k,1c}^*$。例如，参与者 1 的未来个人禀赋在三个状态都是零，他仍会在总禀赋较多的状态消费得更多。

    \item 所求份额是个人在状态 $\omega$ 的消费占该状态总消费的比例。在均衡状态价格 $\phi_\omega^*$ 下，(a) 给出个人最优消费，(b) 给出总消费与总禀赋的关系：
    \[
       c_{k,1\omega}^*=\frac{c_{k,0}^*}{\sqrt{3\phi_\omega^*}},
       \qquad
       E_{1\omega}=c_{1,1\omega}^*+c_{2,1\omega}^*
       =\frac{c_{1,0}^*+c_{2,0}^*}{\sqrt{3\phi_\omega^*}}
       =\frac{100}{\sqrt{3\phi_\omega^*}}.
    \]

    将两式分别代入分子和分母，约去共同的因子 $1/\sqrt{3\phi_\omega^*}$，再代入 (b) 表中两人的当期消费，得到
    \[
       \frac{c_{k,1\omega}^*}{E_{1\omega}}
       =\frac{c_{k,0}^*/\sqrt{3\phi_\omega^*}}{100/\sqrt{3\phi_\omega^*}}
       =\frac{c_{k,0}^*}{100}
       =\begin{cases}
           (600/13)/100=6/13, & k=1,\\
           (700/13)/100=7/13, & k=2.
         \end{cases}
    \]

    这两个份额分别等于两人在 0 期的消费份额，在 $a,b,c$ 三个状态都相同。由 (c) 的比例式 $c_{k,1\omega}^*=(c_{k,0}^*/100)E_{1\omega}$，代入 $E_{1a}=200$、$E_{1c}=50$，得到
    \[
       \frac{c_{k,1a}^*}{c_{k,1c}^*}
       =\frac{(c_{k,0}^*/100)\cdot 200}{(c_{k,0}^*/100)\cdot 50}
       =\frac{200}{50}=4,\qquad k=1,2.
    \]

    从状态 $c$ 到状态 $a$，总禀赋和两人的消费量都变为原来的四倍。两人以各自的固定份额分担总禀赋的波动。
\end{enumerate}
\end{solution}

\endgroup
